% The triangle of orbit 21 with the local flag complex above it: the figure of
% figures/triangle_cell.tex with, on its left, the incidence graph of
% figures/triangle_incidence.tex and a two-headed arrow onto the cell for the
% projection, a surjection. Hand-placed through the hook \gtwocelloverlay of
% julia/to_tex/triangle_cell_figure.jl, in the coordinates of the triangle
% (generated with --scale=0.5 --taxon-scale=0.5 --ray-scale=0.8 --name-below
% --no-cell-letters --cell-overhang --height=4.3, the height of the square of
% figures/square_lettered.tex): its centre at the origin, its top corner at
% (0,2.8667), its top circle of radius 1.2032, its left side from (0,2.8667)
% to (-2.4826,-1.4333). The graph is on the left, at 0.64 of its size as in
% figures/triangle_projection_5.tex, the middle of its height (-0.691 in its
% own coordinates) at that of the picture of the triangle (from -2.650 to
% 4.419, so 0.8846), and its right end 0.25 short of the arrow; the arrow, 2.4
% long, is horizontal at that height and stops 0.15 short of the left side,
% at (-1.144,0.8846), outside the triangle. The whole figure is drawn at 0.95.
\begingroup
\def\gtwocelloverlay{%
  \begin{scope}[shift={(-6.725,1.576)},scale=.64,transform shape]
    % The local flag complex above the triangle of orbit 21 (type \mathfrak A_1),
% as the incidence graph of its tetrahedra and of the triangles that are
% common facets of several of them, each drawn as such, with its vertices. Hand-drawn; it is the tree
% of the cell with every edge subdivided and its seven leaves deleted: a
% triangle at every inner vertex, a tetrahedron at the middle of every inner
% edge and at the leaf of every pendant edge. The tree is that of
% figures/triangle_cell.tex, its centre v11 at the origin, its edges in the
% directions they have there, the inner ones of length 2.1 and the pendant
% ones 1.15, and its triangles and tetrahedra are those of
% figures/triangle_incidence_5.tex, so that the two graphs have one scale.
% Source: julia/studies/g2/atlas/flag/local_flag_cell_atlas.csv, rows 21 (QQ
% cell 155, rays 7,12,40), with the letters of figures/triangle_cell.tex:
%   R_7^A = A, R_7^B = B, R_12^A = C, R_12^B = D, R_40^A = F, R_40^B = E,
% e_k = taxon k. Tetrahedra: pendant edges 1BCE, 5BCE (vertex v8), 2BDF, 6BDF
% (v9), 3ADE, 4ADE (v10), 7BDE (v11); inner edges BCDE (v8-v11), BDEF
% (v9-v11), ABDE (v10-v11). Triangles, the node rows: BCE (node 8), BDF (9),
% ADE (10), BDE (11), each the common facet of the tetrahedra of the edges at
% its vertex; no other triangle lies in two tetrahedra.
% The letters are placed so that every triangle has its vertices at the same
% corners as the face of each adjacent cell that they span: top, bottom left,
% bottom right for the faces top-left-front, top-front-right and
% top-left-right of a tetrahedron, and left to right on its bottom face
% left-front-right.
% A scope for figures/triangle_projection.tex, drawn inside the picture of
% figures/triangle_cell.tex at the place that the caller shifts it to; its
% names start with i, apart from those of the trees.
\begin{scope}[x=1.35cm,y=1.35cm]
\definecolor{g2TaxonBlue}{HTML}{1676A3}
\definecolor{g2TaxonYellow}{HTML}{E39A24}
\definecolor{g2HeartRed}{HTML}{D7443E}
\definecolor{g2RayLabel}{HTML}{000000}
\definecolor{g2CellLabel}{HTML}{000000}
\tikzset{
  g2inc edge/.style={draw=g2RayLabel,line width=.7pt},
  g2inc label/.style={text=g2RayLabel,inner sep=0pt,font={\sffamily\bfseries\fontsize{7}{7}\selectfont}},
  g2inc solid label/.style={g2inc label,font={\sffamily\bfseries\fontsize{9.5}{9.5}\selectfont}},
  % A triangle: an equilateral one, its vertices at its corners, the top one
  % first, then the left one.
  pics/g2tri/.style n args={3}{code={
    \draw[g2CellLabel,fill=g2TaxonBlue!12,line width=.6pt,line join=round]
      (0mm,6.3mm)--(-5.456mm,-3.15mm)--(5.456mm,-3.15mm)--cycle;
    \node[g2inc solid label] at (0mm,2.12mm) {#1};
    \node[g2inc solid label] at (-2.06mm,-1.47mm) {#2};
    \node[g2inc solid label] at (2.06mm,-1.47mm) {#3};
  }},
  % A tetrahedron: its two front faces, lit from the left, and its vertices
  % at its corners: the top one, the left one, the front one, the right one;
  % a pendant one has its taxon at the top, drawn by the caller.
  pics/g2tet faces/.style={code={
    \path[fill=g2TaxonBlue!8] (0mm,6.27mm)--(-7.48mm,-2.86mm)--(1.1mm,-5.39mm)--cycle;
    \path[fill=g2TaxonBlue!26] (0mm,6.27mm)--(1.1mm,-5.39mm)--(7.48mm,-1.98mm)--cycle;
    \draw[g2CellLabel,line width=.6pt,line join=round] (0mm,6.27mm)--(-7.48mm,-2.86mm)--(1.1mm,-5.39mm)--(7.48mm,-1.98mm)--cycle;
  }},
  pics/g2tet/.style n args={4}{code={
    \pic {g2tet faces};
    \node[g2inc solid label] at (.22mm,3.3mm) {#1};
    \pic {g2tet base={#2}{#3}{#4}};
  }},
  pics/g2tet base/.style n args={3}{code={
    \node[g2inc solid label] at (-4.29mm,-1.98mm) {#1};
    \node[g2inc solid label] at (.99mm,-3.19mm) {#2};
    \node[g2inc solid label] at (4.51mm,-1.43mm) {#3};
  }},
  pics/g2pendant/.style n args={3}{code={
    \pic {g2tet faces};
    \pic {g2tet base={#1}{#2}{#3}};
  }},
  g2inc dark/.style={fill=g2TaxonBlue,draw=white,line width=.4pt,inner sep=0pt},
  g2inc light/.style={fill=g2TaxonYellow,draw=white,line width=.4pt,inner sep=0pt},
}
% The vertices of the G2-tree of the cell, which is not drawn.
\coordinate (iv1) at (2.17001,1.09585);
\coordinate (iv2) at (0.90086,-2.81480);
\coordinate (iv3) at (-2.65243,-1.06025);
\coordinate (iv4) at (-2.17001,1.09585);
\coordinate (iv5) at (2.65243,-1.06025);
\coordinate (iv6) at (-0.90086,-2.81480);
\coordinate (iv7) at (0.00000,1.15000);
\coordinate (iv8) at (2.09936,-0.05197);
\coordinate (iv9) at (0.00000,-2.10000);
\coordinate (iv10) at (-2.09936,-0.05197);
\coordinate (iv11) at (0.00000,0.00000);
% The tetrahedra of the inner edges, at their middles, that of v8-v11 moved
% 0.017 of the edge towards v8 to even its gaps to the two triangles.
\coordinate (ie8) at (barycentric cs:iv8=.517,iv11=.483);
\coordinate (ie9) at (barycentric cs:iv9=1,iv11=1);
\coordinate (ie10) at (barycentric cs:iv10=1,iv11=1);
% A tetrahedron and the triangles among its facets.
\foreach \a/\b in {iv1/iv8,iv5/iv8,iv2/iv9,iv6/iv9,iv3/iv10,iv4/iv10,iv7/iv11,
                   ie8/iv8,ie8/iv11,ie9/iv9,ie9/iv11,ie10/iv10,ie10/iv11}
  \draw[g2inc edge] (\a)--(\b);
\pic at (iv8) {g2tri={E}{B}{C}};
\pic at (iv9) {g2tri={F}{D}{B}};
\pic at (iv10) {g2tri={E}{A}{D}};
\pic at (iv11) {g2tri={E}{D}{B}};
\pic at (iv1) {g2pendant={B}{E}{C}};
\pic at (iv2) {g2pendant={D}{F}{B}};
\pic at (iv3) {g2pendant={A}{E}{D}};
\pic at (iv4) {g2pendant={A}{E}{D}};
\pic at (iv5) {g2pendant={B}{E}{C}};
\pic at (iv6) {g2pendant={D}{F}{B}};
\pic at (iv7) {g2pendant={D}{E}{B}};
\pic at (ie8) {g2tet={E}{D}{B}{C}};
\pic at (ie9) {g2tet={E}{D}{F}{B}};
\pic at (ie10) {g2tet={E}{A}{D}{B}};
% The taxa of the pendant tetrahedra, as in the trees, at their top corners.
\node[g2inc dark,rectangle,minimum size=3.45mm] at ([xshift=0mm,yshift=2.25mm]iv1) {};
\begin{scope}[shift={([xshift=0mm,yshift=2.25mm]iv2)},scale=.675]
\path[fill=g2TaxonBlue,draw=white,line width=.4pt,line join=round] (0mm,3.0mm)--(-3.0mm,-2.2mm)--(3.0mm,-2.2mm)--cycle;
\end{scope}
\node[g2inc dark,circle,minimum size=3.75mm] at ([xshift=0mm,yshift=2.25mm]iv3) {};
\node[g2inc light,rectangle,minimum size=3.45mm] at ([xshift=0mm,yshift=2.25mm]iv4) {};
\begin{scope}[shift={([xshift=0mm,yshift=2.25mm]iv5)},scale=.675]
\path[fill=g2TaxonYellow,draw=white,line width=.4pt,line join=round] (0mm,3.0mm)--(-3.0mm,-2.2mm)--(3.0mm,-2.2mm)--cycle;
\end{scope}
\node[g2inc light,circle,minimum size=3.75mm] at ([xshift=0mm,yshift=2.25mm]iv6) {};
\begin{scope}[shift={([xshift=0mm,yshift=2.25mm]iv7)},scale=.675]
\path[fill=g2HeartRed,draw=white,line width=.4pt]
  (0mm,-2.7mm)..controls(-.45mm,-2.1mm)and(-3.0mm,.35mm)..(-3.0mm,1.45mm)
  ..controls(-3.0mm,3.2mm)and(-.9mm,3.4mm)..(0mm,1.75mm)
  ..controls(.9mm,3.4mm)and(3.0mm,3.2mm)..(3.0mm,1.45mm)
  ..controls(3.0mm,.35mm)and(.45mm,-2.1mm)..(0mm,-2.7mm)--cycle;
\end{scope}
\end{scope}

  \end{scope}
  \draw[g2RayLabel,line width=.9pt,
    -{Stealth[length=2.2mm,width=2mm]Stealth[length=2.2mm,width=2mm]}]
    (-3.694,0.8846)--(-1.294,0.8846);
}%
\scalebox{.95}{% Coarsest G2 Grassmannian cell 155: rays 7,12,40, orbit 21, with its rays clockwise from the top corner; G2-trees with lettered inner vertices, generated by julia/to_tex/triangle_cell_figure.jl
% Bounded edges of length 1.6 in the cell and 2.0 in the rays, before scaling.
% Letters of the inner vertices of the rays, in the order of the rays:
%   A: the side {3,4} of its edges
%   B: the side {1,2,5,6,7} of its edges
%   C: the side {1,5} of its edges
%   D: the side {2,3,4,6,7} of its edges
%   E: the side {1,3,4,5,7} of its edges
%   F: the side {2,6} of its edges
\begingroup
\definecolor{g2RayLabel}{HTML}{000000}
\definecolor{g2CellLabel}{HTML}{000000}
\definecolor{g2TaxonBlue}{HTML}{1676A3}
\colorlet{g2SquareFill}{g2TaxonBlue!10}
\tikzset{
g2square side/.style={draw=g2RayLabel,line width=.8pt},
g2square circle/.style={draw=g2RayLabel,fill=white,line width=.8pt},
g2square cell/.style={fill=g2SquareFill},
g2square name/.style={text=g2RayLabel,inner sep=0pt,font=\boldmath},
g2square cell name/.style={g2square name,text=g2CellLabel},
  g2triangle letter/.style={text=g2RayLabel,inner sep=0pt,font={\sffamily\bfseries\fontsize{8}{8}\selectfont}},
}
\begin{tikzpicture}[x=1cm,y=1cm,baseline=(current bounding box.center)]
\fill[g2square cell] (0.0,2.8667)--(2.4826,-1.4333)--(-2.4826,-1.4333)--cycle;
\draw[g2square side] (0.0,2.8667)--(2.4826,-1.4333)--(-2.4826,-1.4333)--cycle;
% Ray 7
\draw[g2square circle] (0.0,2.8667) circle[radius=1.2032];
\node[g2square name,anchor=south] at (0.0,4.1499) {\textpdfrender{TextRenderingMode=FillStroke,LineWidth=.5pt}{\bfseries\boldmath \(\alpha_3\)}};
\begin{scope}[shift={(0.0,2.8667)}]
\begingroup
\definecolor{g2TaxonBlue}{HTML}{1676A3}
\definecolor{g2TaxonYellow}{HTML}{E39A24}
\definecolor{g2HeartRed}{HTML}{D7443E}
\tikzset{
  g2tree pendant/.style={draw=black,line width=.6pt,line cap=round},
  g2tree internal/.style={draw=black,line width=2pt,line cap=round},
  g2tree vertex/.style={circle,fill=black,draw=white,line width=.4pt,inner sep=0pt,minimum size=3.8pt},
  g2tree dark/.style={fill=g2TaxonBlue,draw=white,line width=.55pt,text=white,inner sep=0pt},
  g2tree light/.style={fill=g2TaxonYellow,draw=white,line width=.55pt,text=white,inner sep=0pt},
  g2tree weight/.style={circle,fill=white,draw=none,text=black,inner sep=.15pt,minimum size=8.5pt},
  g2tree aligned pendant/.style={fill=white,text=black,inner sep=.7pt,sloped,above=1pt,font={\sffamily\fontsize{7.5}{8}\selectfont}},
  g2tree aligned internal/.style={fill=white,text=black,inner sep=1pt,sloped,above=1pt,font={\sffamily\fontsize{7.5}{8}\selectfont\bfseries}},
}
\coordinate (v1) at (0.66299,0.52872);
\coordinate (v2) at (0.36793,-0.76402);
\coordinate (v3) at (-0.82674,-0.18870);
\coordinate (v4) at (-0.66299,0.52872);
\coordinate (v5) at (0.82674,-0.18870);
\coordinate (v6) at (-0.36793,-0.76402);
\coordinate (v7) at (0.00000,0.84800);
\coordinate (v8) at (-0.38997,0.08901);
\coordinate (v9) at (0.38997,-0.08901);
\draw[g2tree internal] (v8)--(v9);
\draw[g2tree pendant] (v1)--(v9);
\draw[g2tree pendant] (v2)--(v9);
\draw[g2tree pendant] (v3)--(v8);
\draw[g2tree pendant] (v4)--(v8);
\draw[g2tree pendant] (v5)--(v9);
\draw[g2tree pendant] (v6)--(v9);
\draw[g2tree pendant] (v7)--(v9);
\node[g2tree vertex] at (v8) {};
\node[g2tree vertex] at (v9) {};
\node[g2tree dark,rectangle,minimum size=2.484mm,font={\sffamily\fontsize{4.8}{4.8}\selectfont\bfseries}] at (v1) {};
\begin{scope}[shift={(v2)},scale=0.54,transform shape]
\path[fill=g2TaxonBlue,draw=white,line width=.55pt,line join=round] (0mm,3.0mm)--(-3.0mm,-2.2mm)--(3.0mm,-2.2mm)--cycle;
\end{scope}
\node[g2tree dark,circle,minimum size=2.592mm,font={\sffamily\fontsize{4.8}{4.8}\selectfont\bfseries}] at (v3) {};
\node[g2tree light,rectangle,minimum size=2.484mm,font={\sffamily\fontsize{4.8}{4.8}\selectfont\bfseries}] at (v4) {};
\begin{scope}[shift={(v5)},scale=0.54,transform shape]
\path[fill=g2TaxonYellow,draw=white,line width=.55pt,line join=round] (0mm,3.0mm)--(-3.0mm,-2.2mm)--(3.0mm,-2.2mm)--cycle;
\end{scope}
\node[g2tree light,circle,minimum size=2.592mm,font={\sffamily\fontsize{4.8}{4.8}\selectfont\bfseries}] at (v6) {};
\begin{scope}[shift={(v7)},scale=0.54,transform shape]
\path[fill=g2HeartRed,draw=white,line width=.55pt]
  (0mm,-2.7mm)..controls(-.45mm,-2.1mm)and(-3.0mm,.35mm)..(-3.0mm,1.45mm)
  ..controls(-3.0mm,3.2mm)and(-.9mm,3.4mm)..(0mm,1.75mm)
  ..controls(.9mm,3.4mm)and(3.0mm,3.2mm)..(3.0mm,1.45mm)
  ..controls(3.0mm,.35mm)and(.45mm,-2.1mm)..(0mm,-2.7mm)--cycle;
\end{scope}
\node[g2triangle letter] at (-0.3456,-0.1682) {A};
\node[g2triangle letter] at (0.107,0.1496) {B};
\endgroup
\end{scope}
% Ray 12
\draw[g2square circle] (2.4826,-1.4333) circle[radius=1.2032];
\node[g2square name,anchor=north west] at (3.5939,-2.0749) {\textpdfrender{TextRenderingMode=FillStroke,LineWidth=.5pt}{\bfseries\boldmath \(\alpha_3\)}};
\begin{scope}[shift={(2.4826,-1.4333)}]
\begingroup
\definecolor{g2TaxonBlue}{HTML}{1676A3}
\definecolor{g2TaxonYellow}{HTML}{E39A24}
\definecolor{g2HeartRed}{HTML}{D7443E}
\tikzset{
  g2tree pendant/.style={draw=black,line width=.6pt,line cap=round},
  g2tree internal/.style={draw=black,line width=2pt,line cap=round},
  g2tree vertex/.style={circle,fill=black,draw=white,line width=.4pt,inner sep=0pt,minimum size=3.8pt},
  g2tree dark/.style={fill=g2TaxonBlue,draw=white,line width=.55pt,text=white,inner sep=0pt},
  g2tree light/.style={fill=g2TaxonYellow,draw=white,line width=.55pt,text=white,inner sep=0pt},
  g2tree weight/.style={circle,fill=white,draw=none,text=black,inner sep=.15pt,minimum size=8.5pt},
  g2tree aligned pendant/.style={fill=white,text=black,inner sep=.7pt,sloped,above=1pt,font={\sffamily\fontsize{7.5}{8}\selectfont}},
  g2tree aligned internal/.style={fill=white,text=black,inner sep=1pt,sloped,above=1pt,font={\sffamily\fontsize{7.5}{8}\selectfont\bfseries}},
}
\coordinate (v1) at (0.66299,0.52872);
\coordinate (v2) at (0.36793,-0.76402);
\coordinate (v3) at (-0.82674,-0.18870);
\coordinate (v4) at (-0.66299,0.52872);
\coordinate (v5) at (0.82674,-0.18870);
\coordinate (v6) at (-0.36793,-0.76402);
\coordinate (v7) at (0.00000,0.84800);
\coordinate (v8) at (0.38997,0.08901);
\coordinate (v9) at (-0.38997,-0.08901);
\draw[g2tree internal] (v8)--(v9);
\draw[g2tree pendant] (v1)--(v8);
\draw[g2tree pendant] (v2)--(v9);
\draw[g2tree pendant] (v3)--(v9);
\draw[g2tree pendant] (v4)--(v9);
\draw[g2tree pendant] (v5)--(v8);
\draw[g2tree pendant] (v6)--(v9);
\draw[g2tree pendant] (v7)--(v9);
\node[g2tree vertex] at (v8) {};
\node[g2tree vertex] at (v9) {};
\node[g2tree dark,rectangle,minimum size=2.484mm,font={\sffamily\fontsize{4.8}{4.8}\selectfont\bfseries}] at (v1) {};
\begin{scope}[shift={(v2)},scale=0.54,transform shape]
\path[fill=g2TaxonBlue,draw=white,line width=.55pt,line join=round] (0mm,3.0mm)--(-3.0mm,-2.2mm)--(3.0mm,-2.2mm)--cycle;
\end{scope}
\node[g2tree dark,circle,minimum size=2.592mm,font={\sffamily\fontsize{4.8}{4.8}\selectfont\bfseries}] at (v3) {};
\node[g2tree light,rectangle,minimum size=2.484mm,font={\sffamily\fontsize{4.8}{4.8}\selectfont\bfseries}] at (v4) {};
\begin{scope}[shift={(v5)},scale=0.54,transform shape]
\path[fill=g2TaxonYellow,draw=white,line width=.55pt,line join=round] (0mm,3.0mm)--(-3.0mm,-2.2mm)--(3.0mm,-2.2mm)--cycle;
\end{scope}
\node[g2tree light,circle,minimum size=2.592mm,font={\sffamily\fontsize{4.8}{4.8}\selectfont\bfseries}] at (v6) {};
\begin{scope}[shift={(v7)},scale=0.54,transform shape]
\path[fill=g2HeartRed,draw=white,line width=.55pt]
  (0mm,-2.7mm)..controls(-.45mm,-2.1mm)and(-3.0mm,.35mm)..(-3.0mm,1.45mm)
  ..controls(-3.0mm,3.2mm)and(-.9mm,3.4mm)..(0mm,1.75mm)
  ..controls(.9mm,3.4mm)and(3.0mm,3.2mm)..(3.0mm,1.45mm)
  ..controls(3.0mm,.35mm)and(.45mm,-2.1mm)..(0mm,-2.7mm)--cycle;
\end{scope}
\node[g2triangle letter] at (0.3456,-0.1682) {C};
\node[g2triangle letter] at (-0.0315,-0.1812) {D};
\endgroup
\end{scope}
% Ray 40
\draw[g2square circle] (-2.4826,-1.4333) circle[radius=1.2032];
\node[g2square name,anchor=north east] at (-3.5939,-2.0749) {\textpdfrender{TextRenderingMode=FillStroke,LineWidth=.5pt}{\bfseries\boldmath \(\alpha_3\)}};
\begin{scope}[shift={(-2.4826,-1.4333)}]
\begingroup
\definecolor{g2TaxonBlue}{HTML}{1676A3}
\definecolor{g2TaxonYellow}{HTML}{E39A24}
\definecolor{g2HeartRed}{HTML}{D7443E}
\tikzset{
  g2tree pendant/.style={draw=black,line width=.6pt,line cap=round},
  g2tree internal/.style={draw=black,line width=2pt,line cap=round},
  g2tree vertex/.style={circle,fill=black,draw=white,line width=.4pt,inner sep=0pt,minimum size=3.8pt},
  g2tree dark/.style={fill=g2TaxonBlue,draw=white,line width=.55pt,text=white,inner sep=0pt},
  g2tree light/.style={fill=g2TaxonYellow,draw=white,line width=.55pt,text=white,inner sep=0pt},
  g2tree weight/.style={circle,fill=white,draw=none,text=black,inner sep=.15pt,minimum size=8.5pt},
  g2tree aligned pendant/.style={fill=white,text=black,inner sep=.7pt,sloped,above=1pt,font={\sffamily\fontsize{7.5}{8}\selectfont}},
  g2tree aligned internal/.style={fill=white,text=black,inner sep=1pt,sloped,above=1pt,font={\sffamily\fontsize{7.5}{8}\selectfont\bfseries}},
}
\coordinate (v1) at (0.66299,0.52872);
\coordinate (v2) at (0.36793,-0.76402);
\coordinate (v3) at (-0.82674,-0.18870);
\coordinate (v4) at (-0.66299,0.52872);
\coordinate (v5) at (0.82674,-0.18870);
\coordinate (v6) at (-0.36793,-0.76402);
\coordinate (v7) at (0.00000,0.84800);
\coordinate (v8) at (0.00000,-0.40000);
\coordinate (v9) at (-0.00000,0.40000);
\draw[g2tree internal] (v8)--(v9);
\draw[g2tree pendant] (v1)--(v9);
\draw[g2tree pendant] (v2)--(v8);
\draw[g2tree pendant] (v3)--(v9);
\draw[g2tree pendant] (v4)--(v9);
\draw[g2tree pendant] (v5)--(v9);
\draw[g2tree pendant] (v6)--(v8);
\draw[g2tree pendant] (v7)--(v9);
\node[g2tree vertex] at (v8) {};
\node[g2tree vertex] at (v9) {};
\node[g2tree dark,rectangle,minimum size=2.484mm,font={\sffamily\fontsize{4.8}{4.8}\selectfont\bfseries}] at (v1) {};
\begin{scope}[shift={(v2)},scale=0.54,transform shape]
\path[fill=g2TaxonBlue,draw=white,line width=.55pt,line join=round] (0mm,3.0mm)--(-3.0mm,-2.2mm)--(3.0mm,-2.2mm)--cycle;
\end{scope}
\node[g2tree dark,circle,minimum size=2.592mm,font={\sffamily\fontsize{4.8}{4.8}\selectfont\bfseries}] at (v3) {};
\node[g2tree light,rectangle,minimum size=2.484mm,font={\sffamily\fontsize{4.8}{4.8}\selectfont\bfseries}] at (v4) {};
\begin{scope}[shift={(v5)},scale=0.54,transform shape]
\path[fill=g2TaxonYellow,draw=white,line width=.55pt,line join=round] (0mm,3.0mm)--(-3.0mm,-2.2mm)--(3.0mm,-2.2mm)--cycle;
\end{scope}
\node[g2tree light,circle,minimum size=2.592mm,font={\sffamily\fontsize{4.8}{4.8}\selectfont\bfseries}] at (v6) {};
\begin{scope}[shift={(v7)},scale=0.54,transform shape]
\path[fill=g2HeartRed,draw=white,line width=.55pt]
  (0mm,-2.7mm)..controls(-.45mm,-2.1mm)and(-3.0mm,.35mm)..(-3.0mm,1.45mm)
  ..controls(-3.0mm,3.2mm)and(-.9mm,3.4mm)..(0mm,1.75mm)
  ..controls(.9mm,3.4mm)and(3.0mm,3.2mm)..(3.0mm,1.45mm)
  ..controls(3.0mm,.35mm)and(.45mm,-2.1mm)..(0mm,-2.7mm)--cycle;
\end{scope}
\node[g2triangle letter] at (0.0,-0.661) {F};
\node[g2triangle letter] at (-0.1696,0.071) {E};
\endgroup
\end{scope}
% The cell
\node[g2square cell name,anchor=north] at (0.0,-1.5133) {\textpdfrender{TextRenderingMode=FillStroke,LineWidth=.5pt}{\bfseries\boldmath \oldstylenums{21}}};
\begingroup
\definecolor{g2TaxonBlue}{HTML}{1676A3}
\definecolor{g2TaxonYellow}{HTML}{E39A24}
\definecolor{g2HeartRed}{HTML}{D7443E}
\tikzset{
  g2tree pendant/.style={draw=black,line width=.6pt,line cap=round},
  g2tree internal/.style={draw=black,line width=2pt,line cap=round},
  g2tree vertex/.style={circle,fill=black,draw=white,line width=.4pt,inner sep=0pt,minimum size=3.8pt},
  g2tree dark/.style={fill=g2TaxonBlue,draw=white,line width=.55pt,text=white,inner sep=0pt},
  g2tree light/.style={fill=g2TaxonYellow,draw=white,line width=.55pt,text=white,inner sep=0pt},
  g2tree weight/.style={circle,fill=white,draw=none,text=black,inner sep=.15pt,minimum size=8.5pt},
  g2tree aligned pendant/.style={fill=white,text=black,inner sep=.7pt,sloped,above=1pt,font={\sffamily\fontsize{7.5}{8}\selectfont}},
  g2tree aligned internal/.style={fill=white,text=black,inner sep=1pt,sloped,above=1pt,font={\sffamily\fontsize{7.5}{8}\selectfont\bfseries}},
}
\coordinate (v1) at (0.82874,0.66090);
\coordinate (v2) at (0.45992,-0.95503);
\coordinate (v3) at (-1.03342,-0.23587);
\coordinate (v4) at (-0.82874,0.66090);
\coordinate (v5) at (1.03342,-0.23587);
\coordinate (v6) at (-0.45992,-0.95503);
\coordinate (v7) at (0.00000,1.06000);
\coordinate (v8) at (0.79976,0.19010);
\coordinate (v9) at (0.00000,-0.59010);
\coordinate (v10) at (-0.79976,0.19010);
\coordinate (v11) at (0.00000,0.20990);
\draw[g2tree internal] (v8)--(v11);
\draw[g2tree internal] (v9)--(v11);
\draw[g2tree internal] (v10)--(v11);
\draw[g2tree pendant] (v1)--(v8);
\draw[g2tree pendant] (v2)--(v9);
\draw[g2tree pendant] (v3)--(v10);
\draw[g2tree pendant] (v4)--(v10);
\draw[g2tree pendant] (v5)--(v8);
\draw[g2tree pendant] (v6)--(v9);
\draw[g2tree pendant] (v7)--(v11);
\node[g2tree vertex] at (v8) {};
\node[g2tree vertex] at (v9) {};
\node[g2tree vertex] at (v10) {};
\node[g2tree vertex] at (v11) {};
\node[g2tree dark,rectangle,minimum size=3.105mm,font={\sffamily\fontsize{6}{6}\selectfont\bfseries}] at (v1) {};
\begin{scope}[shift={(v2)},scale=0.675,transform shape]
\path[fill=g2TaxonBlue,draw=white,line width=.55pt,line join=round] (0mm,3.0mm)--(-3.0mm,-2.2mm)--(3.0mm,-2.2mm)--cycle;
\end{scope}
\node[g2tree dark,circle,minimum size=3.24mm,font={\sffamily\fontsize{6}{6}\selectfont\bfseries}] at (v3) {};
\node[g2tree light,rectangle,minimum size=3.105mm,font={\sffamily\fontsize{6}{6}\selectfont\bfseries}] at (v4) {};
\begin{scope}[shift={(v5)},scale=0.675,transform shape]
\path[fill=g2TaxonYellow,draw=white,line width=.55pt,line join=round] (0mm,3.0mm)--(-3.0mm,-2.2mm)--(3.0mm,-2.2mm)--cycle;
\end{scope}
\node[g2tree light,circle,minimum size=3.24mm,font={\sffamily\fontsize{6}{6}\selectfont\bfseries}] at (v6) {};
\begin{scope}[shift={(v7)},scale=0.675,transform shape]
\path[fill=g2HeartRed,draw=white,line width=.55pt]
  (0mm,-2.7mm)..controls(-.45mm,-2.1mm)and(-3.0mm,.35mm)..(-3.0mm,1.45mm)
  ..controls(-3.0mm,3.2mm)and(-.9mm,3.4mm)..(0mm,1.75mm)
  ..controls(.9mm,3.4mm)and(3.0mm,3.2mm)..(3.0mm,1.45mm)
  ..controls(3.0mm,.35mm)and(.45mm,-2.1mm)..(0mm,-2.7mm)--cycle;
\end{scope}
\endgroup
\ifdefined\gtwocelloverlay\gtwocelloverlay\fi
\end{tikzpicture}
\endgroup
}%
\endgroup
